\chapter*{译后记}
\addcontentsline{toc}{chapter}{译后记} % 手动加入目录
\markboth{译后记}{译后记} % 设置页眉

本书作者为当代德国著名哲学家曼弗雷德·里德尔(Manfred Riedel，1936—2009)。里德尔于1936年5月出生在莱比锡的一个乡村，1954年进入莱比锡大学，1957年转入海德堡大学，曾先后师从恩斯特·布洛赫、卡尔·洛维特、伽达默尔等当代德国哲学大师。1960年，他在洛维特的指导下以《黑格尔思想中的理论与实践》一文获得博士学位。1968年，他完成了教授资格论文《市民社会：一个古典政治学和现代自然法范畴》。他曾先后在海德堡大学、马堡大学、萨尔布吕肯大学任教，1970年成为埃朗根-纽伦堡大学哲学正教授，1980—1981年担任纽约社会研究新学院休斯(Theodor-Heuss)客座教授，并先后在都灵、罗马、威尼斯等地担任客座教授。1992—1993年，他担任耶拿大学席勒讲座教授，1992年获得哈勒-维滕堡大学的实践哲学教席，并于2004年荣休。此外，里德尔曾于1991年至2003年间担任德国海德格尔协会主席，并从2005年起担任佛罗伦萨意大利人文科学院院士。里德尔的研究涉及康德、黑格尔、尼采、狄尔泰和海德格尔等人，从古代、中世纪到欧洲近代的市民社会史，构成其研究的一个重点，他的《在传统与革命之间：黑格尔法哲学研究》即为这方面研究的一部力作，也是当代德国学界黑格尔法哲学研究方面的一部经典作品。

该书第1版于1969年以《黑格尔法哲学研究》(Studien zu Hegels Rechtsphilosophie)之名出版(Frankfurt am Main: Suhrkamp, 1969)，次年刊行第2版。1982年，里德尔将原书内容稍加扩充，并加上主标题“在传统与革命之间”（原标题则保留为副标题）后，由位于斯图加特的克莱特-科塔(Klett-Cotta)出版社出版，是为第3版。自第1版出版以来，先后有意大利文、日文、西班牙文、韩文、英文等诸种译本问世，其影响也随之扩展到国际学界。当原著出版50周年之际，其中译本能够面世，固然有些恨晚，但也不能不说是一件幸事。

本书的翻译出版首先要感谢商务印书馆“政治哲学名著译丛”主编吴彦博士的盛情邀请。此次翻译根据的是里德尔原书1982年扩充版，并且参考了赖特(W. Wright)的英译本(Manfred Riedel, Between Tradition and Revolution: The Hegelian Transformation of Political Philosophy, Cambridge University Press 1984)。翻译的具体分工是：前言部分由两位译者共同完成，第一、第二部分的翻译由黄钰洲承担，并由朱学平进行了校对；第三、第四部分的翻译由朱学平承担。“术语索引”和“人名索引”由黄钰洲录入。需要指出的是，我国罗马法学者普遍将“jus civile”译为“市民法”，本书也同样将“societas civilis”译为“市民社会”，其含义是“公民社会”，而非黑格尔意义上的“市民社会”。由于译者学养和语言水平有限，译文中必然存在错误不当之处，对此，译者真诚期待读者和学界同人不吝批评指正。最后，我们也衷心感谢商务印书馆的编校者为本书出版所付出的辛勤劳动！

\begin{flushright}
译者\\
2019年1月
\end{flushright}